% 第一节：矢量与对偶矢量
\chapter{矢量与对偶矢量}

\section{矢量的坐标无关描述}

在中学的立体几何中，一个矢量常被写成有序数组 $(x,y,z)$。这种说法既有道理，又不准确。
准确的地方在于，对任一个三维矢量空间总可以建立坐标系，从而把矢量写成有序数组；不准确的地方
在于，这样的坐标系有无数种，没有哪一种比别的更特殊。既然没有理由选取某一个坐标系，最好的
做法就是一个也不选；一个也不选，就不能把矢量写成有序数组。

矢量和有序数组在概念上是有区别的：有序数组是在矢量上附加了一个结构（一个指定的坐标系）之后
得到的东西。物理中有类似的情形：就基本物理规律而言一切惯性系平权，然而在地面上讨论问题时
地面参考系是特殊的，这时便相当于给物理附加了一个结构。

那么什么是矢量？粗略地说，抽象的矢量就是去掉了坐标系这一附加结构的有序数组。精确的说法如下。

\begin{definition}[矢量空间]
设 $V$ 是一个集合，在 $V$ 上定义了加法 $V\times V\to V$ 与数乘 $\R\times V\to V$，它们满足
交换律、结合律、分配律以及存在零元、存在负元等通常的运算律。若存在 $V$ 中的 $n$ 个元素
$e_1,\dots,e_n$，使得 $V$ 的每个元素都能唯一地写成它们的线性组合，则称 $V$ 为 $n$ 维矢量空间，
称其元素为矢量，称 $\{e_1,\dots,e_n\}$ 为 $V$ 的一组基。
\end{definition}

定义中断言了基的存在，却没有指定基是哪一个。因此，在抽象的矢量空间中除了加法与数乘之外
并没有别的东西，特别地，无法自然地定义点乘。点乘是由坐标系定义的，因而无法用它来判断两个
矢量是否垂直。能够只用加法与数乘定义的概念是共线与共面。

\begin{definition}[共线与共面]
若存在不全为零的实数 $\alpha,\beta$ 使得 $\alpha A+\beta B=0$，称矢量 $A,B$ 共线；若存在
不全为零的实数 $\alpha,\beta,\gamma$ 使得 $\alpha A+\beta B+\gamma C=0$，称矢量 $A,B,C$ 共面。
\end{definition}

\section{对偶空间与对偶基}

\begin{definition}[线性函数与对偶空间]
映射 $f:V\to\R$ 称为 $V$ 上的函数。若对任意 $A,B\in V$ 与任意实数 $\alpha,\beta$ 恒有
\[
f(\alpha A+\beta B)=\alpha f(A)+\beta f(B),
\]
则称 $f$ 为 $V$ 上的线性函数，简称线性函数。$V$ 上全体线性函数构成的集合记为 $V^*$，
在其中定义加法与数乘
\[
(f+g)(A)=f(A)+g(A),\qquad (\alpha f)(A)=\alpha f(A).
\]
在这两种运算下 $V^*$ 构成矢量空间，称为 $V$ 的对偶空间。
\end{definition}

线性函数的运算律是逐点定义的，因此交换律、分配律等自动成立，$V^*$ 确实构成矢量空间。

\begin{proposition}[对偶基]
设 $V$ 是 $n$ 维矢量空间，$\{e_1,\dots,e_n\}$ 是 $V$ 的一组基。则存在 $V^*$ 中唯一的一组元素
$\{e^1,\dots,e^n\}$ 满足
\[
e^i(e_j)=\delta^i_j=
\begin{cases}
1,& i=j,\\
0,& i\neq j .
\end{cases}
\]
这组元素构成 $V^*$ 的一组基，称为 $\{e_i\}$ 的对偶基。
\end{proposition}

\begin{proof}
设 $A=\sum_i A^i e_i$。对给定的线性函数 $f$，由线性性
\[
f(A)=\sum_i A^i f(e_i),
\]
故 $f$ 由 $n$ 个实数 $f(e_1),\dots,f(e_n)$ 唯一确定；反过来任给 $n$ 个实数都可以这样定义一个线性函数。
所以 $\dim V^*=n$。再验证 $\{e^i\}$ 的线性无关性与生成性：若 $\sum_i c_i e^i=0$，把它作用于 $e_j$
得 $c_j=0$，故线性无关；又对任意 $f$，记 $f_i=f(e_i)$，则 $f$ 与 $\sum_i f_i e^i$ 在基矢 $e_j$ 上
取值相同，从而作为线性函数相等。
\end{proof}

把 $f(e_i)$ 简记为 $f_i$，则任一 $f\in V^*$ 可唯一地写成
\[
f=\sum_i f_i e^i,\qquad f(A)=\sum_i f_iA^i .
\]

\section{对偶对与点的表示}

上面构造中 $V$ 与 $V^*$ 的地位并不对称：一个是矢量空间，一个是它上面的线性函数空间。但函数
$f$ 作用于 $A$ 得到 $f(A)$，也可以理解为把 $A$ 看作定义在 $V^*$ 上的函数，作用于自变量 $f$ 得到
同一个数。既然两者的地位可以互换，就把 $f(A)$ 写成
\[
f\cdot A=f_iA^i .
\]
这一运算很像点乘，但与点乘有一个本质区别：

\begin{remark}
矢量与矢量之间、对偶矢量与对偶矢量之间没有自然的缩并，只有矢量与对偶矢量之间才有。因此
形如 $A^iB^i$ 或 $A_iB_i$ 的记号没有意义，而 $A^iB_i$ 有意义。
\end{remark}

这一区别正是后文把矢量解释为点、把对偶矢量解释为直线的起点：点是矢量，直线是对偶矢量，
“点与直线有包含关系”这一几何事实将表现为它们之间的缩并。
